% Generated draft: requires mathematical and source review.
\documentclass[11pt]{article}
\usepackage[margin=25mm]{geometry}
\usepackage{fontspec}
\setmainfont{Latin Modern Roman}
\usepackage{amsmath,amssymb,mathtools,unicode-math}
\setmathfont{Latin Modern Math}
\usepackage{enumitem,needspace,xurl,hyperref}
\setlength{\parindent}{0pt}
\setlength{\parskip}{5pt}
\setlength{\emergencystretch}{3em}
\newcommand{\N}{\mathbb{N}}
\newcommand{\Z}{\mathbb{Z}}
\newcommand{\Q}{\mathbb{Q}}
\newcommand{\R}{\mathbb{R}}
\newcommand{\C}{\mathbb{C}}
\newcommand{\F}{\mathbb{F}}
\newcommand{\A}{\mathbb{A}}
\newcommand{\PP}{\mathbb{P}}
\newcommand{\E}{\mathbb{E}}
\newcommand{\eps}{\varepsilon}
\newcommand{\dd}{\,\mathrm{d}}
\begin{document}
\section*{UnsolvedMath ID: 30003618}
\textbf{Source problem code:} catalogue entry
\subsection{Definitions and mathematical statement}
The supplied record does not specify a mathematical definition of a quasicrystal, a class of tilings, or a construction and topology for the associated tiling spaces. Accordingly, these terms retain their source-dependent meanings; no particular model or additional hypotheses are imposed here.

\textbf{Methodological question.} What algebro-topological tools should be employed to study mathematical quasicrystals and their tiling spaces?

This is a broad request to identify suitable methods, rather than a proposition with specified hypotheses and conclusion. The source does not specify the properties to be studied or criteria for the suitability of a tool. A more precise quantified formulation therefore requires additional information about the objects, their tiling spaces, and the intended mathematical objectives.
\subsection{Short English statement}
The question asks which algebraic and topological methods are suitable for studying mathematical quasicrystals and their tiling spaces. The supplied record leaves the class of objects, the construction of the spaces, and the goals of the study unspecified.
\subsection{Sources}
\begin{itemize}
\item Supplied local UnsolvedMath registry record, permanent ID \(30003618\), title \emph{Algebraic Topology of Mathematical Quasicrystals}, source problem code \texttt{catalogue entry}, in the collection \emph{Oberwolfach Reports Open Problems}. The supplied proposed year is 2017; no individual proposer is identified.
\item Supplied report citation: \emph{Spectral Structures and Topological Methods in Mathematical Quasicrystals} (2018), Oberwolfach Reports, pp.~2781--2845, DOI \texttt{10.4171/owr/2017/46}, \url{https://doi.org/10.4171/owr/2017/46}. The same link is supplied under literature sources with the label \emph{Oberwolfach Report 46 2017}. The background describes these reports as documenting mathematical workshops and their open-problem discussions.
\item Literature assessment dated exactly \texttt{2026-08-21}: the supplied curation labels the question open and describes it as a broad methodological question about suitable algebro-topological tools, rather than a proposition with a known proof or disproof. This is supplied dated curation, not an independently verified assessment. Its stated checked source is the report linked above. The supplied background expressly requests independent verification of time-sensitive status and cited claims.
\end{itemize}
\end{document}
